\subsection{Equações do Movimento de Atitude nos referenciais LVLH e ECI}
\label{sec:atitude_inercial}

A formulação matemática da atitude da espaçonave perseguidora é obtida a partir das equações de Euler aplicadas a um corpo rígido, escritas em relação aos eixos fixos no corpo da espaçonave.
A atitude pode ser representada em diferentes referenciais por meio de matrizes de rotação: no sistema inercial (ECI), para uma descrição absoluta, ou no sistema orbital \gls{lvlh}, para uma descrição relativa à órbita.
Em todos os casos, a atitude é definida em coordenadas do corpo e transformada para o referencial escolhido por meio dessas matrizes.

Com base em \cite{wie2008space}, a equação do movimento angular de um corpo rígido em torno do seu \gls{cdm} pode ser escrita, no referencial inercial, como

\begin{equation}
\vec{M}
=
\dot{\vec{H}},
\label{eq:wie_M_Hdot}
\end{equation}

em que $\vec{H}$ é a quantidade de movimento angular e $\vec{M}$ é o momento externo resultante.

\begin{figure}[H]
\centering
\begin{tikzpicture}[scale=1.0, line cap=round, line join=round]
  % Origem e deslocamento entre referenciais
  \coordinate (O_N) at (0,0);
  \coordinate (O_B) at (6,0);

  % Referencial N (inercial) -- geometria aproximada à Fig. 1.2 de Wie
  \draw[->, thick] (O_N) -- ++(-1.4,-1.4) node[below left] {$\vec{n}_x$};
  \draw[->, thick] (O_N) -- ++(2.6,0)     node[below right] {$\vec{n}_y$};
  \draw[->, thick] (O_N) -- ++(0,2.4)     node[above] {$\vec{n}_z$};
  \node[below] at ($(O_N)+(0.1,-0.2)$) {$N$};

  % Referencial B (corpo) -- eixos inclinados (mesmo corpo, base diferente)
  \draw[->, thick] (O_B) -- ++(-0.9,-2.1) node[below left] {$\vec{b}_x$};
  \draw[->, thick] (O_B) -- ++(2.3,0.4)   node[below right] {$\vec{b}_y$};
  \draw[->, thick] (O_B) -- ++(-0.6,2.2)  node[above left] {$\vec{b}_z$};
  \node[below] at ($(O_B)+(0.1,-0.2)$) {$B$};

  % Velocidade angular (B em relacao a N)
  \draw[->, very thick] ($(O_B)+(0.15,0.15)$) -- ++(2.0,2.0)
    node[above right] {$\vec{\omega}^{B/N}$};
\end{tikzpicture}
\caption{Dois sistemas de coordenadas ($N$ e $B$) associados ao mesmo corpo, com $B$ girando em relação a $N$ com velocidade angular $\vec{\omega}^{B/N}$ (adaptado de \cite{wie2008space}).}
\label{fig:wie_ref_N_B}
\end{figure}

A Figura~\ref{fig:wie_ref_N_B} ilustra dois referenciais associados ao mesmo corpo: $N$ (inercial) e $B$ (fixo ao corpo), em que $B$ gira em relação a $N$ com velocidade angular $\vec{\omega}^{B/N}$.

Aplicando o teorema do transporte entre os referenciais $N$ e $B$, para um vetor $\vec{H}$ expresso em $B$, obtém-se

\begin{equation}
\dot{\vec{H}}
\equiv
\left\{ \frac{d\vec{H}}{dt} \right\}_{N}
=
\left\{ \frac{d\vec{H}}{dt} \right\}_{B}
+
\vec{\omega}^{B/N} \times \vec{H},
\label{eq:wie_transport_H}
\end{equation}

onde $\vec{\omega}^{B/N}$ denota a velocidade angular de $B$ em relação a $N$, expressa no referencial $B$.
%%%%%%%%%%%%%%%%%%%%%%%%%%%%%%%%%%%%%%%%%%%%%%%%%%%%%%%%
Para um corpo rígido com eixos principais fixos ao corpo, no referencial $B$ a quantidade de movimento angular é

\begin{equation}
\vec{H}
= \mathbf{I}\,\vec{\omega}^{B/N},
\label{eq:wie_H_Jw}
\end{equation}

onde $\mathbf{I}$ é o tensor de inércia expresso em $B$, constante nesse referencial.
Adotando a notação $(\dot{\,})$ para derivadas temporais tomadas no referencial do corpo $B$ e identificando $\vec{\omega}\equiv\vec{\omega}^{B/N}$, obtém-se

\begin{equation}
\left\{ \frac{d\vec{H}}{dt} \right\}_{B}
=
\mathbf{I}\,\dot{\vec{\omega}}^{B/N},
\label{eq:wie_dHdt_B}
\end{equation}

pois $\mathbf{I}$ não depende do tempo em $B$.


%%%%%%%%%%%%%%%%%%%%%%%%%%%%%%%%%%%%%%%%%%%%%%%%%%%%%%%%
Substituindo \eqref{eq:wie_dHdt_B} e \eqref{eq:wie_H_Jw} em \eqref{eq:wie_transport_H} e, em seguida, em \eqref{eq:wie_M_Hdot}, obtém-se

\begin{equation}
\vec{M}
=
\mathbf{I}\,\dot{\vec{\omega}}^{B/N}
+
\vec{\omega}^{B/N} \times
\big( \mathbf{I}\,\vec{\omega}^{B/N} \big),
\label{eq:euler_rotacional_wie}
\end{equation}

que é a forma vetorial das equações de Euler para o movimento de rotação de um corpo rígido em torno do seu \gls{cdm}, com grandezas expressas no referencial do corpo.

Identificando $\vec{\omega} \equiv \vec{\omega}^{B/N}$, obtém-se
\begin{equation}
\mathbf{I}\,\dot{\vec{\omega}}
+
\vec{\omega}\times\big(\mathbf{I}\,\vec{\omega}\big)
=
\vec{M}.
\label{eq:euler_corpo_rigido}
\end{equation}
em que $\mathbf{I}$ é o tensor de inércia em relação ao \gls{cdm}, $\vec{\omega}$ é a velocidade angular e $\vec{M}$ é o momento externo resultante aplicado ao corpo.

%%%%%%%%%%%%%%%%%%%%%%%%%%%%%%%%%%%%%%%%%%%%%%%%%%%%%%%%
A equação \eqref{eq:euler_corpo_rigido} é válida para a espaçonave perseguidora na fase não operacional do \gls{mr}.
Essa expressão é obtida como caso particular do modelo com inércia variável \cite{wie2008space}, quando o \gls{mr} está travado e, portanto, $\dot{\mathbf{I}}_B=\mathbf{0}$ e $\mathbf{I}_B$ pode ser considerada constante.
A partir de \eqref{eq:euler_corpo_rigido}, obtêm-se as equações escalares assumindo que o sistema de coordenadas do corpo está alinhado com os eixos principais de inércia, de modo que

\begin{equation}
   \mathbf{I} =
\begin{bmatrix}
I_x & 0   & 0 \\
0   & I_y & 0 \\
0   & 0   & I_z
\end{bmatrix},
\qquad
\vec{\omega} =
\begin{bmatrix}
\omega_x \\[2pt]
\omega_y \\[2pt]
\omega_z
\end{bmatrix},
\qquad
\vec{M} =
\begin{bmatrix}
M_x \\[2pt]
M_y \\[2pt]
M_z
\end{bmatrix}. 
\end{equation}

%%%%%%%%%%%%%%%%%%%%%%%%%%%%%%%%%%%%%%%%%%%%%%%%%%%%%%%%

Primeiro, calcula-se

\begin{equation}
\mathbf{I}\,\vec{\omega}
=
\begin{bmatrix}
I_x \omega_x \\
I_y \omega_y \\
I_z \omega_z
\end{bmatrix}.
\label{eq:Iomega_componentes}
\end{equation}

Em seguida, calcula-se o produto vetorial, em que $\{\vec{b}_1,\vec{b}_2,\vec{b}_3\}$ é a base (ortonormal) do referencial do corpo $B$:

\begin{equation}
\vec{\omega} \times (\mathbf{I}\vec{\omega})
=
\begin{vmatrix}
\vec{b}_1 & \vec{b}_2 & \vec{b}_3 \\
\omega_x & \omega_y & \omega_z \\
I_x\omega_x & I_y\omega_y & I_z\omega_z
\end{vmatrix}
\end{equation}

cujo resultado é

\begin{equation}
\vec{\omega} \times (\mathbf{I}\vec{\omega})
=
\begin{bmatrix}
(I_z - I_y)\,\omega_y\,\omega_z \\
(I_x - I_z)\,\omega_x\,\omega_z \\
(I_y - I_x)\,\omega_x\,\omega_y
\end{bmatrix}.
\label{eq:cross_Iomega}
\end{equation}

%%%%%%%%%%%%%%%%%%%%%%%%%%%%%%%%%%%%%%%%%%%%%%%%%%%%%%%%

A derivada do termo inercial é dada por

\begin{equation}
\mathbf{I}\,\dot{\vec{\omega}}
=
\begin{bmatrix}
I_x \dot{\omega}_x \\
I_y \dot{\omega}_y \\
I_z \dot{\omega}_z
\end{bmatrix}.
\label{eq:I_omega_dot}
\end{equation}

Substituindo \eqref{eq:cross_Iomega} e \eqref{eq:I_omega_dot} em \eqref{eq:euler_corpo_rigido} e igualando componente a componente, obtêm-se as equações de Euler na forma escalar:
\begin{equation}
\begin{aligned}
I_x \dot{\omega}_x + (I_z - I_y)\,\omega_y\,\omega_z &= M_x, \\
I_y \dot{\omega}_y + (I_x - I_z)\,\omega_x\,\omega_z &= M_y, \\
I_z \dot{\omega}_z + (I_y - I_x)\,\omega_x\,\omega_y &= M_z.
\end{aligned}
\end{equation}

Essas três equações diferenciais ordinárias são acopladas e não lineares.
Além disso, descrevem a evolução das componentes da velocidade angular do corpo rígido.
Em conjunto com as equações cinemáticas correspondentes, essas equações definem o movimento rotacional de um corpo com três graus de liberdade rotacionais \cite{wie2008space}.

%%%%%%%%%%%%%%%%%%%%%%%%%%%%%%%%%%%%%%%%%%%%%%%%%%%%%%%%
\subsubsection{Matriz de inércia}
\label{sec:tensor_inercia_variavel}

No caso da espaçonave perseguidora, a distribuição de massa do sistema depende da configuração do \gls{mr}.
Para descrever a atitude da espaçonave, adotam-se os parâmetros de Euler (quatérnions) $\vec{q}_B(t)$.
Seja definido o vetor de coordenadas do sistema:

\begin{equation}
\vec{\xi}(t) =
\begin{bmatrix}
\vec{q}_B(t) \\
\vec{q}_m(t)
\end{bmatrix},
\qquad
\vec{q}_m(t) =
\begin{bmatrix}
q_1(t) \\ q_2(t) \\ q_3(t)
\end{bmatrix},
\qquad
\vec{q}_B^{\,T}(t)\,\vec{q}_B(t)=1,
\end{equation}

onde $\vec{q}_B(t)$ parametriza a atitude da espaçonave e $\vec{q}_m(t)$ contém os
ângulos das três juntas revolutas do \gls{mr}.
Para simplificar a notação, nesta subseção escreve-se $\vec{q}(t)\equiv\vec{q}_m(t)$.
A matriz de inércia total em torno do centro de massa, expressa no referencial do
corpo $B$, é modelada por

\begin{equation}
\mathbf{I}_B\big(\vec{q}(t)\big)
=
\mathbf{I}_{\text{perseg},B}
+
\mathbf{I}_{\text{manip},B}\big(\vec{q}(t)\big),
\label{eq:inercia_total}
\end{equation}

onde $\mathbf{I}_{\text{perseg},B}$ é a matriz de inércia da espaçonave isolada e
$\mathbf{I}_{\text{manip},B}$ representa a contribuição do \gls{mr}, dependente
da configuração articular e da orientação relativa dos elos em relação à espaçonave.


A equação fundamental da quantidade de movimento angular de um corpo rígido em
torno de seu centro de massa, apresentada em Wie~\cite{wie2008space}, relaciona
o momento externo $\vec{M}$ à derivada temporal da quantidade de movimento
angular $\vec{H}$ tomada no referencial inercial $N$:

\begin{equation}
\vec{M}
=
\left\{ \frac{d\vec{H}}{dt} \right\}_{N}.
\label{eq:eq_momento_angular}
\end{equation}

Para o sistema em estudo, o vetor da quantidade de movimento angular expresso no
referencial do corpo é

\begin{equation}
\vec{H}_B(t)
=
\mathbf{I}_B\big(\vec{q}(t)\big)\,
\vec{\omega}^{B/N}(t),
\label{eq:H_def}
\end{equation}
em que $\vec{\omega}^{B/N}$ é a velocidade angular da espaçonave em relação a $N$.

Aplicando o teorema de transporte entre os referenciais $N$ e $B$, conforme \cite{wie2008space}, obtém-se
\begin{equation}
\left\{ \frac{d\vec{H}_B}{dt} \right\}_{N}
=
\left\{ \frac{d\vec{H}_B}{dt} \right\}_{B}
+
\vec{\omega}^{B/N} \times \vec{H}_B.
\label{eq:teorema_transporte_H}
\end{equation}

Substituindo \eqref{eq:H_def} em \eqref{eq:teorema_transporte_H} e usando
\eqref{eq:eq_momento_angular}, obtém-se a equação de movimento rotacional para o sistema
com inércia variável:

\begin{equation}
\vec{M}
=
\mathbf{I}_B\big(\vec{q}(t)\big)\,
\dot{\vec{\omega}}^{B/N}
+
\dot{\mathbf{I}}_B\big(\vec{q}(t)\big)\,
\vec{\omega}^{B/N}
+
\vec{\omega}^{B/N} \times
\Big(
\mathbf{I}_B\big(\vec{q}(t)\big)\,
\vec{\omega}^{B/N}
\Big),
\label{eq:euler_inercia_variavel}
\end{equation}
em que o termo
$\vec{\omega}^{B/N} \times \big(\mathbf{I}_B(\vec{q}(t))\,\vec{\omega}^{B/N}\big)$
corresponde ao efeito giroscópico clássico das equações de Euler, e o termo
$\dot{\mathbf{I}}_B(\vec{q}(t))\,\vec{\omega}^{B/N}$ representa a contribuição
associada à variação temporal da matriz de inércia devido à redistribuição
interna de massas provocada pelo movimento do \gls{mr} sobre a espaçonave.

% \begin{equation}
%   \begin{bmatrix}
%     a & b & c & d & e\\
%     1 &  &  & 4 & 5
%   \end{bmatrix}
% \end{equation}